\documentclass[12pt]{ximera}
\title{Final Exam (Spring 2023)}
\begin{document}
\begin{abstract}
An archived cumulative final: true/false across all chapters, Borda count and Plurality-with-Elimination, Banzhaf and Shapley-Shubik power, lone-divider and sealed bids, and Hamilton, Huntington-Hill, and Webster apportionment.
\end{abstract}
\maketitle

\section*{Final Exam (Spring 2023)}

\begin{problem}
True or false?

\begin{prompt}
\textbf{1.} A candidate who wins a majority of first-place votes also has a plurality.
\begin{multipleChoice}
\choice[correct]{True}
\choice{False}
\end{multipleChoice}
\begin{feedback}
A majority is more than everyone else combined, so certainly the most.
\end{feedback}
\end{prompt}

\begin{prompt}
\textbf{2.} A Condorcet candidate wins all head-to-head matchups against the other candidates.
\begin{multipleChoice}
\choice[correct]{True}
\choice{False}
\end{multipleChoice}
\begin{feedback}
That is the definition.
\end{feedback}
\end{prompt}

\begin{prompt}
\textbf{3.} It is possible to invent a voting method that satisfies all the fairness criteria.
\begin{multipleChoice}
\choice{True}
\choice[correct]{False}
\end{multipleChoice}
\begin{feedback}
Arrow's Impossibility Theorem says no method can satisfy all of them in every election.
\end{feedback}
\end{prompt}

\begin{prompt}
\textbf{4.} A dictator has enough weight to pass any motion in a weighted voting system.
\begin{multipleChoice}
\choice[correct]{True}
\choice{False}
\end{multipleChoice}
\begin{feedback}
A dictator's weight alone meets the quota.
\end{feedback}
\end{prompt}

\begin{prompt}
\textbf{5.} In a sequential coalition, the order in which players join does not matter.
\begin{multipleChoice}
\choice{True}
\choice[correct]{False}
\end{multipleChoice}
\begin{feedback}
Order is the whole point: it determines the pivotal player.
\end{feedback}
\end{prompt}

\begin{prompt}
\textbf{6.} A fair share depends on the number of players in a sharing problem.
\begin{multipleChoice}
\choice[correct]{True}
\choice{False}
\end{multipleChoice}
\begin{feedback}
With $N$ players, a fair share is worth at least $\frac1N$ of the total.
\end{feedback}
\end{prompt}

\begin{prompt}
\textbf{7.} In the lone-divider method, it is impossible for all choosers to bid for the same share.
\begin{multipleChoice}
\choice{True}
\choice[correct]{False}
\end{multipleChoice}
\begin{feedback}
All choosers may find only one and the same share acceptable. That is a standoff, which the method resolves.
\end{feedback}
\end{prompt}

\begin{prompt}
\textbf{8.} The geometric mean of two numbers $a$ and $b$ is $\sqrt{a+b}$.
\begin{multipleChoice}
\choice{True}
\choice[correct]{False}
\end{multipleChoice}
\begin{feedback}
It is $\sqrt{ab}$.
\end{feedback}
\end{prompt}

\begin{prompt}
\textbf{9.} Some apportionment methods violate the quota rule.
\begin{multipleChoice}
\choice[correct]{True}
\choice{False}
\end{multipleChoice}
\begin{feedback}
Jefferson's, Adams', Webster's, and Huntington-Hill's methods all can.
\end{feedback}
\end{prompt}

\begin{prompt}
\textbf{10.} With Hamilton's method, adding one seat to the total can leave some state with fewer seats than before.
\begin{multipleChoice}
\choice[correct]{True}
\choice{False}
\end{multipleChoice}
\begin{feedback}
That is the Alabama paradox.
\end{feedback}
\end{prompt}
\end{problem}

\begin{problem}
A quiz bowl team heading home from a tournament ranks three fast-food stops: Arby's (A),
Burger King (B), and Culver's (C).
\begin{center}
\begin{tabular}{c|cccc}
Number of voters & 12 & 8 & 5 & 2\\\hline
1st & A & C & B & B\\
2nd & B & B & A & C\\
3rd & C & A & C & A
\end{tabular}
\end{center}
Use the Borda count. A $\answer{56}$ \quad B $\answer{61}$ \quad C $\answer{45}$ \qquad
Winner: $\answer{B}$

\begin{feedback}
A: $12\cdot3+8\cdot1+5\cdot2+2\cdot1=56$; B: $12\cdot2+8\cdot2+5\cdot3+2\cdot3=61$;
C: $12\cdot1+8\cdot3+5\cdot1+2\cdot2=45$.
\end{feedback}
\end{problem}

\begin{problem}
Aberforth (A), Bellatrix (B), Cedric (C), and Diagon (D) are running for Minister of
Magic.
\begin{center}
\begin{tabular}{c|ccccc}
Number of voters & 20 & 15 & 8 & 5 & 2\\\hline
1st & A & B & D & C & C\\
2nd & B & D & C & B & A\\
3rd & C & A & B & A & D\\
4th & D & C & A & D & B
\end{tabular}
\end{center}
Use Plurality-with-Elimination. Eliminated first: $\answer{C}$ \quad second:
$\answer{D}$ \qquad Winner: $\answer{B}$

\begin{feedback}
Round 1: A $20$, B $15$, D $8$, C $7$. C's $5$ votes go to B and $2$ go to A. Round 2:
A $22$, B $20$, D $8$. D's $8$ votes go to B (C is gone). Round 3: B $28$, A $22$.
\end{feedback}
\end{problem}

\begin{problem}
Consider the weighted voting system $[18:10,8,5,2]$.

\begin{prompt}
\textbf{(a)} How many winning coalitions are there? $\answer{4}$
\end{prompt}

\begin{prompt}
\textbf{(b)} Banzhaf power distribution: $\beta_1=\answer{1/2}$, $\beta_2=\answer{1/2}$,
$\beta_3=\answer{0}$, $\beta_4=\answer{0}$
\begin{feedback}
Winning: $\set{P_1,P_2}$, $\set{P_1,P_2,P_3}$, $\set{P_1,P_2,P_4}$, and the grand
coalition. Only $P_1$ and $P_2$ are ever critical ($4$ times each); $P_3$ and $P_4$ are
dummies.
\end{feedback}
\end{prompt}
\end{problem}

\begin{problem}
Consider the weighted voting system $[15:10,8,5]$. Find the Shapley-Shubik power
distribution.

$\sigma_1=\answer{2/3}$ \quad $\sigma_2=\answer{1/6}$ \quad $\sigma_3=\answer{1/6}$

\begin{feedback}
$\seq{P_1,\underline{P_2},P_3}$, $\seq{P_1,\underline{P_3},P_2}$,
$\seq{P_2,\underline{P_1},P_3}$, $\seq{P_2,P_3,\underline{P_1}}$,
$\seq{P_3,\underline{P_1},P_2}$, $\seq{P_3,P_2,\underline{P_1}}$.
\end{feedback}
\end{problem}

\begin{problem}
Four friends divide a cake by the lone-divider method. Don is the divider.
\begin{center}
\begin{tabular}{r|cccc}
 & $s_1$ & $s_2$ & $s_3$ & $s_4$\\\hline
Don & 25\% & 25\% & 25\% & 25\%\\
Connie & 80\% & 0\% & 10\% & 10\%\\
Cecelia & 5\% & 5\% & 40\% & 50\%\\
Caleb & 30\% & 10\% & 30\% & 30\%
\end{tabular}
\end{center}

\begin{prompt}
\textbf{(a)} Which shares does Caleb bid for?
\begin{selectAll}
\choice[correct]{$s_1$}
\choice{$s_2$}
\choice[correct]{$s_3$}
\choice[correct]{$s_4$}
\end{selectAll}
\begin{feedback}
Bids (at least $25\%$): Connie $\set{s_1}$, Cecelia $\set{s_3,s_4}$, Caleb
$\set{s_1,s_3,s_4}$.
\end{feedback}
\end{prompt}

\begin{prompt}
\textbf{(b)} How many fair distributions are there? $\answer{2}$ \qquad Don always gets
share number $\answer{2}$.
\begin{feedback}
Connie must take $s_1$, so Cecelia and Caleb split $s_3$ and $s_4$ in either order, and
Don gets $s_2$.
\end{feedback}
\end{prompt}
\end{problem}

\begin{problem}
Three fans sneak backstage after a concert and find an earring and a used napkin that
they believe belonged to the famous singer. They divide them by sealed bids.
\begin{center}
\begin{tabular}{r|rrr}
 & Alex & Bill & Christie\\\hline
Earring & \$500 & \$400 & \$300\\
Napkin & \$100 & \$200 & \$150
\end{tabular}
\end{center}

\begin{prompt}
\textbf{(a)} Fair shares: Alex $\$\answer{200}$, Bill $\$\answer{200}$, Christie
$\$\answer{150}$
\end{prompt}

\begin{prompt}
\textbf{(b)} Initial settlement (positive = pays): Alex $\answer{300}$, Bill $\answer{0}$,
Christie $\answer{-150}$
\begin{feedback}
Alex gets the earring ($500$, which is $300$ over), Bill the napkin (exactly his fair
share), and Christie nothing.
\end{feedback}
\end{prompt}

\begin{prompt}
\textbf{(c)} Final settlement (positive = pays): Alex $\answer{250}$, Bill $\answer{-50}$,
Christie $\answer{-200}$
\begin{feedback}
The surplus is $300-150=150$, or $\$50$ each.
\end{feedback}
\end{prompt}
\end{problem}

\begin{problem}
The Lylat system apportions $435$ senate seats among four planets by Hamilton's method.
\begin{center}
\begin{tabular}{r|cccc}
Planet & Fortuna & Katina & Corneria & Sauria\\\hline
Population & 104,240 & 100,960 & 200,450 & 29,350
\end{tabular}
\end{center}
Apportionment: Fortuna $\answer{104}$, Katina $\answer{101}$, Corneria $\answer{201}$,
Sauria $\answer{29}$

\begin{feedback}
$SD=435{,}000/435=1000$, so the quotas are $104.24$, $100.96$, $200.45$, $29.35$. The
lower quotas total $433$; the extra $2$ go to Katina ($0.96$) and Corneria ($0.45$).
\end{feedback}
\end{problem}

\begin{problem}
Five colleges at a university form a council of $60$ seats: Agriculture (A), Business
(B), Design (D), Engineering (E), and Graduate (G).
\begin{center}
\begin{tabular}{r|ccccc}
College & A & B & D & E & G\\\hline
Population & 971 & 1,363 & 148 & 3,171 & 347
\end{tabular}
\end{center}

\begin{prompt}
\textbf{(a)} Standard divisor: $\answer{100}$
\end{prompt}

\begin{prompt}
\textbf{(b)} Huntington-Hill (modified quotas to $3$ decimals): which divisor from this
list works?
\begin{multipleChoice}
\choice{$100$}
\choice[correct]{$101$}
\choice{$99$}
\end{multipleChoice}

\smallskip
Apportionment: A $\answer{10}$, B $\answer{14}$, D $\answer{2}$, E $\answer{31}$,
G $\answer{3}$
\begin{feedback}
At $SD=100$ the quotas $9.71$, $13.63$, $1.48$, $31.71$, $3.47$ all exceed their cutoffs
($\sqrt{90}\approx9.487$, $\sqrt{182}\approx13.491$, $\sqrt2\approx1.414$,
$\sqrt{992}\approx31.496$, $\sqrt{12}\approx3.464$), giving $62$. At $d=101$ the quotas
are $9.614$, $13.495$, $1.465$, $31.396$, $3.436$; E and G drop below their cutoffs,
for a total of $60$.
\end{feedback}
\end{prompt}

\begin{prompt}
\textbf{(c)} Webster at the standard divisor: A $\answer{10}$, B $\answer{14}$,
D $\answer{1}$, E $\answer{32}$, G $\answer{3}$

\smallskip
Is it the same as Huntington-Hill?
\begin{multipleChoice}
\choice{Yes}
\choice[correct]{No}
\end{multipleChoice}
\begin{feedback}
Conventional rounding already totals $60$. Huntington-Hill moves a seat from the largest
college, Engineering, to the smallest, Design.
\end{feedback}
\end{prompt}
\end{problem}

\end{document}
